%!TEX root = main.tex 

\section{Related Work}
\label{sec:related}

\noindent We discuss related work other than those covered in \cref{sec:motivation}.

\noindent\textbf{GPU simulation.}
GPU computation simulators like SCALE-sim~\cite{samajdar2018scale} and Accel-Sim~\cite{khairy2020accel} focus on instruction-level modeling, while other approaches~\cite{hu2022dpro,lu2023distsim,luo2022srifty,zhu2020daydream} rely on profiling or tracing to capture accurate \op execution times. \sys similarly adopts a trace-based methodology, but crucially does so without requiring a full-scale deployment.




\noindent\textbf{Network simulation.} 
% Tools like ASTRA-sim~\cite{won2023astra}, FlexFlow~\cite{jia2019beyond}, Proteus~\cite{duan2024proteus}, and NCCL-Predictor~\cite{nccl} rely on the ${\alpha-\beta}$ model~\cite{valiant1990bridging}, offering simplicity but often limited accuracy. 
Discrete event simulators like ns-3~\cite{riley2010ns} and OMNeT++~\cite{varga2019practical} are well-known to have high overheads due to their packet-level whole-stack simulation. Other than parallelization optimizations \cite{GCLL23,BZTW24}, recently learning based approach to predict packet-level or flow-level performance without simulating the whole stack~\cite{ZNKY21,li2024m3,YPCL22} has shown promising results.
As mentioned before, integrating them with \sys for training simulation is a promising direction for future work. 

\begin{comment}
\noindent\textbf{Concurrent execution.} 
Concurrent GPU execution is increasingly employed to improve utilization~\cite{jiang2024megascale,chen2024centauri,zhang2017poseidon}, but it often degrades computation kernel performance, as these kernels overlap primarily with memcpy or communication operations. While this slowdown is significant (\cref{subsec:challenges}), most simulators ignore it. Proteus~\cite{duan2024proteus} is the only known simulator that attempts to model such slowdown, but its heuristic factor—based solely on GPU architecture and model type—is too coarse to handle diverse, unseen workloads. In contrast, \sys uses an ML-based approach for more accurate and generalizable slowdown predictions.
\end{comment}
